\documentclass[aspectratio=169]{beamer}

% ---------- ENCODING & FONTS ----------
\usepackage[utf8]{inputenc}
\usepackage[T1]{fontenc}

% ---------- THEME & LOOK ----------
\usetheme{Madrid}
\usecolortheme{seahorse}
\usefonttheme{professionalfonts}
\setbeamertemplate{navigation symbols}{}
\setbeamercovered{transparent}
\setbeamertemplate{itemize items}[circle]
\setbeamertemplate{enumerate items}[default]
\setlength{\parskip}{0.25em}

% ---------- COLORS ----------
\usepackage{xcolor}
\definecolor{accent}{RGB}{0,92,184}
\definecolor{soft}{RGB}{233,244,255}
\setbeamercolor{structure}{fg=accent}
\setbeamercolor{title}{fg=white,bg=accent}
\setbeamercolor{frametitle}{fg=accent}
\setbeamercolor{block title}{bg=accent!12,fg=accent}
\setbeamercolor{block body}{bg=soft}

% ---------- PACKAGES ----------
\usepackage{booktabs}
\usepackage{amsmath, amssymb}
\usepackage{array}
\usepackage{multirow}
\usepackage{hyperref}
\hypersetup{colorlinks=true,linkcolor=accent,urlcolor=accent}

% ---------- SAFE FALLBACK FOR TRANSITIONS ----------
\makeatletter
\@ifundefined{transfade}{\newcommand{\transfade}[1][]{}}{}
\makeatother

% ---------- PLACEHOLDER BOX (compile-safe; no images required) ----------
\newcommand{\placeholderbox}[2][0.82\linewidth]{%
  \begin{center}
    \fbox{\begin{minipage}[c][0.46\textheight][c]{#1}
      \centering \vspace{0.25em}
      \textbf{#2}\\[0.6em]
      \emph{Insert figure/diagram here later.}
    \end{minipage}}
  \end{center}
}

% ---------- TITLE ----------
\title{\textbf{Chapter 10: Wage Structures Across Markets  \vspace{.5cm}\\  }}
\subtitle{ {\large ECON 300 - Labour Economics}}
%\institute{UNBC}
\author{Muhebullah Karimzada}

\date{Winter 2026}


\begin{document}

% ============================================================
% TITLE
% ============================================================
\begin{frame}[plain]
  \titlepage
\end{frame}

% ============================================================
\section{Learning Objectives \& Main Questions}
% ============================================================

\begin{frame}{Learning Objectives}
\transfade
\begin{itemize}[<+->]
  \item Understand the Roy model with heterogeneous skills and market sorting.
  \item Explain wage differences across labour markets via supply, demand, and institutions.
  \item Compare wages across occupations, industries, regions, firm size, and sectors.
  \item Interpret interindustry differentials through the lens of efficiency wages.
  \item Discuss why the public sector often pays more than the private sector.
\end{itemize}
\end{frame}

\begin{frame}{Main Questions}
\transfade
\begin{itemize}[<+->]
  \item Why do wages vary across labour markets even for apparently identical workers?
  \item Are public sector workers overpaid?
  \item Why are wages higher in Alberta than in Ontario?
  \item What drives interprovincial migration decisions?
  \item Do higher wages at large firms contradict the neoclassical model?
  \item Which industries/occupations pay the most—and is it due to worker skills?
\end{itemize}
\end{frame}

% ============================================================
\section{Explorations with the 2016 Census}
% ============================================================

\begin{frame}{Earnings Model: Returns to Traits}
\transfade
\small
Individual (log) wages depend on schooling, experience, ability/luck, and other traits:
\[
\ln Y_i = \alpha + r S_i + \beta_1 \mathrm{EXP}_i + \beta_2 \mathrm{EXP}_i^2 + \gamma X_i + \varepsilon_i
\]
\begin{itemize}[<+->]
  \item $Y$: earnings;\; $S$: years of schooling;\; $r$: return to schooling (IRR interpretation);
  \item $\mathrm{EXP}$: experience proxy;\; $X$: controls;\; $\varepsilon$: unobservables.
\end{itemize}
\end{frame}

\begin{frame}{Pure Regional, Occupational, and Industry Premia}
\transfade
\begin{itemize}[<+->]
  \item Conditional wage differences remain after adjusting for $S$, $\mathrm{EXP}$, $X$.
  \item “Premiums” summarize province, occupation, and industry effects.
\end{itemize}
\end{frame}

\begin{frame}{Table 10.1 — Earnings Differentials by Province, Canada, 2015}
\transfade
\scriptsize
\begin{tabular}{lcc}
\toprule
\textbf{Province} & \textbf{\% Jobs} & \textbf{Premium} \\
\midrule
Newfoundland and Labrador & 1.3 & 4.1 \\
Prince Edward Island & 0.3 & $-12.8$ \\
Nova Scotia & 2.7 & $-8.1$ \\
New Brunswick & 2.1 & $-11.1$ \\
Quebec & 22.8 & $-11.5$ \\
Ontario & 39.9 & --- \\
Manitoba & 3.5 & $-2.6$ \\
Saskatchewan & 2.8 & 10.1 \\
Alberta & 13.1 & 18.1 \\
British Columbia & 12.1 & 1.6 \\
\bottomrule
\end{tabular}
\end{frame}

\begin{frame}{Table 10.1 — Earnings Differentials by Occupation and Industry, Canada, 2015}
\transfade
\scriptsize
\begin{tabular}{lcc|lcc}
\toprule
\multicolumn{3}{c}{\textbf{Occupation}} & \multicolumn{3}{c}{\textbf{Industry}} \\
\textbf{Group} & \textbf{\% Jobs} & \textbf{Premium} & \textbf{Group} & \textbf{\% Jobs} & \textbf{Premium} \\
\midrule
Senior management & 2.0 & 56.1 & Agriculture, forestry, and fishing & 1.5 & $-35.1$ \\
Middle management & 5.0 & 28.5 & Other primary & 2.4 & 51.3 \\
Other management & 9.1 & 9.5 & Utilities & 1.5 & 38.5 \\
Professionals (business/finance) & 3.4 & 14.8 & Construction & 10.2 & $-4.2$ \\
Finance and administration & 5.8 & $-15.6$ & Manufacturing & 16.2 & --- \\
Clerical & 1.1 & $-25.6$ & Wholesale trade & 6.4 & 1.7 \\
Professionals (science) & 8.0 & 14.2 & Retail trade & 9.6 & $-23.6$ \\
Other science occupations & 5.8 & $-2.1$ & Transportation/warehousing & 7.1 & $-0.7$ \\
Professionals (health) & 1.0 & 30.9 & Information and culture & 2.9 & 5.4 \\
Other health occupations & 1.0 & 0.9 & Finance and insurance & 4.5 & 16.2 \\
Professionals (education/social) & 4.8 & 10.4 & Real estate and leasing & 1.3 & $-8.9$ \\
Other education/social & 3.0 & 10.0 & Professional services & 7.9 & $-3.1$ \\
Professionals (arts/culture/rec) & 1.4 & $-21.3$ & Admin/support, waste mgmt & 3.5 & $-25.7$ \\
Sales and service supervisors & 4.4 & $-9.5$ & Educational services & 4.5 & $-16.4$ \\
Sales representatives & 5.7 & $-18.7$ & Health and social assistance & 3.5 & $-22.8$ \\
Sales and services support & 4.2 & $-33.4$ & Accommodation and food & 3.6 & $-57.6$ \\
Trade workers & 15.3 & --- & Other services (excl.\ public admin) & 4.6 & $-27.8$ \\
Transportation/trade helpers & 10.1 & $-22.7$ & Public administration & 8.8 & 6.8 \\
Supervisors/tech (primary/secondary) & 5.8 & $-8.2$ &  &  &  \\
Labourers & 3.3 & $-30.4$ &  &  &  \\
\bottomrule
\end{tabular}
\end{frame}

% ============================================================
\section{Occupational Wage Structures}
% ============================================================

\begin{frame}{NOC and Occupational Wage Structure}
\transfade
\begin{itemize}[<+->]
  \item National Occupational Classification (NOC): 26 two-digit major groups; 17 broad classes; hundreds of detailed codes.
  \item Occupational wage gaps reflect nonpecuniary traits, human capital, endowed skills, demand shifts, and institutions.
\end{itemize}
\end{frame}

\begin{frame}{Figure 10.1 — Occupational Wage Differentials \textit{(Placeholder)}}
\transfade
\placeholderbox{Figure 10.1 — Comparative wage levels by occupation.}
\end{frame}

\begin{frame}{Why Do Occupations Pay Differently?}
\transfade
\begin{itemize}[<+->]
  \item \textbf{Compensating differentials}: risk, conditions, schedule.
  \item \textbf{Human capital}: training requirements and entry costs.
  \item \textbf{Endowments}: innate abilities valued differently across jobs.
  \item \textbf{Short-run demand}: cyclical/sectoral shifts.
  \item \textbf{Noncompetitive factors}: licensing, entry limits, unions, wage laws.
\end{itemize}
\end{frame}

% ============================================================
\section{Heterogeneous Workers and the Roy Model}
% ============================================================

\begin{frame}{The Roy Model: Sorting by Comparative Advantage}
\transfade
\begin{itemize}[<+->]
  \item Workers have different skill vectors; wages differ by match quality across occupations.
  \item Individuals self-select into the occupation yielding the highest wage for their skill bundle.
  \item Implication: observed wage distributions reflect \textbf{sorting}, not just treatment effects.
\end{itemize}
\end{frame}

% ============================================================
\section{Regional Wage Structures \& Mobility}
% ============================================================

\begin{frame}{Regional Wage Gaps and Mobility}
\transfade
\begin{itemize}[<+->]
  \item Reasons: geographic preferences; compensating differences (cost of living, remoteness, climate, congestion, pollution).
  \item Short run: disequilibrium encourages mobility; long run: convergence via flows of labour/capital.
  \item Frictions: moving costs, policy barriers (e.g., licensing), and networks.
\end{itemize}
\end{frame}

\begin{frame}{Migration Decision: Benefits vs.\ Costs}
\transfade
\begin{itemize}[<+->]
  \item Move if present value of wage/personal benefits exceeds moving and adaptation costs.
  \item Influences: age, unemployment rates, business cycle, distance, cultural differences.
\end{itemize}
\end{frame}

\begin{frame}{Effect of an Oil Shock on Regional Wages \textit{(Placeholder)}}
\transfade
\placeholderbox{Oil shock — Alberta wage premium; spillovers to construction/services; migration inflows.}
\end{frame}

% ============================================================
\section{Interindustry Wage Differentials}
% ============================================================

\begin{frame}{What Drives Interindustry Wage Gaps?}
\transfade
\begin{itemize}[<+->]
  \item Composition: occupation mix and worker traits.
  \item Region dominance and product-market forces (trade, technology, deregulation).
  \item Nonpecuniary aspects: conditions, seasonality, cyclicality.
  \item Noncompetitive factors: monopoly rents, unions, wage floors, licensing.
\end{itemize}
\end{frame}

\begin{frame}{Efficiency Wages and “Good Jobs”}
\transfade
\begin{itemize}[<+->]
  \item Firms may pay above-market wages to improve morale, deter shirking, reduce turnover, and discourage unionization.
  \item Higher wages can create applicant queues and persistent interindustry premia.
  \item Competitive pressure (imports, deregulation, more rivals) tends to compress premia.
\end{itemize}
\end{frame}

\begin{frame}{Empirical Regularities}
\transfade
\begin{itemize}[<+->]
  \item Interindustry wage structure is \textbf{stable} over time and across countries.
  \item Product-market competition lowers wages relative to protected industries.
  \item Pure interindustry differentials (rents) persist—consistent with efficiency wages.
\end{itemize}
\end{frame}

% ============================================================
\section{Firm Size and Wage Differences}
% ============================================================

\begin{frame}{Interfirm Differentials and Size}
\transfade
\begin{itemize}[<+->]
  \item Larger firms often pay more—due to rents, internal labor markets, scale economies, and sorting.
  \item Also reflect nonpecuniary factors and union presence in larger firms.
\end{itemize}
\end{frame}

% ============================================================
\section{Public vs.\ Private Sector}
% ============================================================

\begin{frame}{Public–Private Wage Differentials}
\transfade
\begin{itemize}[<+->]
  \item Public sector: about one-fifth of Canadian employment; higher wages and benefits on average.
  \item Reasons: job security, benefit generosity, political constraints, monopsony power, unionization, and demand inelasticity.
\end{itemize}
\end{frame}

\begin{frame}{Figure 10.2 — Public Sector Employment, Canada, 1987–2019 \textit{(Placeholder)}}
\transfade
\placeholderbox{Figure 10.2 — Public sector employment share over time.}
\end{frame}

% ============================================================
\section{Chapter Summary}
% ============================================================

\begin{frame}{Summary}
\transfade
\begin{itemize}[<+->]
  \item Wage structure has multiple dimensions: occupation, industry, region, firm size, and sector.
  \item Roy model explains sorting and the mapping from heterogeneous skills to observed wages.
  \item Interindustry and interfirm differentials persist due to rents, institutions, and compensating differences.
  \item Public sector typically pays more; interpretation requires adjusting for traits and benefits.
\end{itemize}
\end{frame}

\begin{frame}
\centering
\Huge{\textbf{Thank You!}} \\

\end{frame}

\end{document}
